\documentclass[UTF8,12pt]{ctexart}
\usepackage[a4paper,margin=2.4cm]{geometry}
\usepackage{graphicx}
\usepackage{booktabs}
\usepackage{longtable}
\usepackage{array}
\usepackage{enumitem}
\usepackage{xcolor}
\usepackage{float}
\usepackage{caption}
\usepackage{hyperref}
\usepackage{tikz}
\usetikzlibrary{arrows.meta,positioning,shapes.geometric,fit,calc}

\hypersetup{
  colorlinks=true,
  linkcolor=blue!55!black,
  urlcolor=blue!55!black
}

\graphicspath{{../screenshots/}}
\setlength{\parindent}{2em}
\setlength{\parskip}{0.35em}
\setlist{nosep,leftmargin=2em}
\renewcommand{\arraystretch}{1.25}
\captionsetup{font=small,labelfont=bf}
\sloppy

\definecolor{mainblue}{HTML}{1F5F99}
\definecolor{lightblue}{HTML}{EAF2FF}
\definecolor{linegray}{HTML}{CED6E0}
\definecolor{softyellow}{HTML}{FFF6DF}

\newcommand{\code}[1]{\texttt{\detokenize{#1}}}
\newcommand{\docmeta}[2]{\noindent\textbf{#1：}#2\par}
\newcommand{\screenshot}[3][0.8\textwidth]{
  \begin{figure}[H]
    \centering
    \includegraphics[width=#1]{#2}
    \caption{#3}
  \end{figure}
}

\tikzset{
  flow/.style={rectangle,rounded corners,draw=mainblue,fill=lightblue,minimum width=2.6cm,minimum height=0.9cm,align=center,font=\small},
  decision/.style={diamond,aspect=2,draw=mainblue,fill=softyellow,align=center,font=\small,inner sep=1pt},
  dataobj/.style={rectangle,draw=linegray,fill=white,minimum width=2.6cm,minimum height=0.85cm,align=center,font=\small},
  actor/.style={rectangle,rounded corners,draw=mainblue,fill=lightblue,minimum width=2.2cm,minimum height=0.9cm,align=center,font=\small\bfseries},
  usecase/.style={ellipse,draw=mainblue,fill=white,minimum width=2.9cm,minimum height=0.85cm,align=center,font=\small},
  classbox/.style={rectangle,draw=mainblue,fill=white,minimum width=3.1cm,minimum height=1.2cm,align=left,font=\small},
  arrow/.style={-{Latex[length=3mm]},thick,draw=mainblue},
  line/.style={thick,draw=mainblue}
}



\title{\textbf{作业管理系统用户使用手册}}
\author{}
\date{}

\begin{document}
\maketitle

\section{启动与登录}

\subsection{Windows 启动方式}

在 Windows 电脑上解压项目后，双击 \code{run_windows.bat}。脚本会自动完成以下步骤：

\begin{enumerate}
  \item 检查 Python 环境。
  \item 创建 \code{.venv} 虚拟环境。
  \item 安装 Flask 依赖。
  \item 检查并初始化 SQLite 数据库。
  \item 自动打开浏览器访问系统。
\end{enumerate}

访问地址：\code{http://127.0.0.1:5000}

\subsection{默认账号}

\begin{table}[H]
\centering
\begin{tabular}{lll}
\toprule
角色 & 用户名 & 密码 \\
\midrule
管理员 & admin & admin123 \\
教师 & teacher01 & teacher123 \\
学生 & student01 & student123 \\
\bottomrule
\end{tabular}
\caption{系统默认演示账号}
\end{table}

\screenshot{01_login.png}{登录页面}

\newpage

\section{管理员使用说明}

管理员登录后进入管理员后台，可以查看用户数、课程数、作业数和提交数。

\screenshot{02_admin_dashboard.png}{管理员后台}

\subsection{教师审核}

\begin{enumerate}
  \item 使用管理员账号登录。
  \item 在“用户管理”表格中找到状态为“待审核”的教师账号。
  \item 点击“通过”按钮。
  \item 教师账号状态变为“正常”后即可登录。
\end{enumerate}

\subsection{停用账号}

\begin{enumerate}
  \item 在用户管理表格中找到需要停用的用户。
  \item 点击“停用”按钮。
  \item 被停用账号再次登录时，系统会提示账号已停用。
\end{enumerate}

\newpage

\section{教师使用说明}

教师登录后进入教师工作台。

\screenshot{03_teacher_dashboard.png}{教师工作台}

\subsection{创建课程}

\begin{enumerate}
  \item 在“开设课程”区域填写课程名称。
  \item 填写课程说明。
  \item 点击“创建课程”。
  \item 创建成功后，课程会显示在“我的课程”列表中。
\end{enumerate}

\subsection{布置作业}

\begin{enumerate}
  \item 在“布置作业”区域选择课程。
  \item 填写作业标题。
  \item 设置截止时间。
  \item 填写作业要求。
  \item 点击“发布作业”。
\end{enumerate}

\subsection{查看提交与评分}

\begin{enumerate}
  \item 在作业列表中点击“查看提交”。
  \item 查看学生提交说明和附件。
  \item 在评分栏填写成绩和评语。
  \item 点击“保存”。
  \item 如需统计，点击“导出统计 CSV”。
\end{enumerate}

\screenshot{07_teacher_submission_detail.png}{教师查看提交页面}

\newpage

\section{学生使用说明}

学生登录后进入学生工作台。

\screenshot{04_student_dashboard.png}{学生工作台}

\subsection{加入课程}

\begin{enumerate}
  \item 在“可选课程”区域查看课程。
  \item 点击“加入课程”。
  \item 加入后课程状态显示为“已加入”。
\end{enumerate}

\subsection{提交作业}

\begin{enumerate}
  \item 在“我的作业”表格中找到需要提交的作业。
  \item 点击“提交/修改”。
  \item 填写提交说明。
  \item 如有附件，点击文件选择框上传。
  \item 点击“保存提交”。
\end{enumerate}

\screenshot{05_student_submit.png}{学生提交作业页面}

提交完成后，学生工作台会显示“已提交”。

\screenshot{06_student_after_submit.png}{提交后的学生工作台}

\subsection{查看成绩}

教师评分后，学生在“我的作业”表格中可以看到成绩和评语。

\end{document}

